\section{Conclusions}
\label{sec:conclusion}

A fluid-filled viscoelastic oblate spheroidal shell model of the human abdomen
reveals two distinct mode families with fundamentally different physics. The
breathing mode ($n = 0$), governed by the fluid bulk modulus, occurs at
approximately \SI{2900}{\hertz} and is irrelevant to infrasonic exposure.
Flexural modes ($n \geq 2$), in which the fluid acts only as added mass, fall
in the \SIrange{4}{10}{\hertz} range for relaxed musculature, consistent with
ISO~2631 whole-body vibration data.

The central finding is a $10^3$--$10^4$ coupling disparity between mechanical
and airborne excitation of these flexural modes. This disparity, arising from
the $(ka)^n$ long-wavelength penalty on acoustic coupling, provides a
quantitative mechanistic explanation for an otherwise puzzling asymmetry in the
empirical literature: mechanical vibration at \SIrange{4}{8}{\hertz}
produces gastrointestinal effects, while airborne infrasound at comparable
frequencies does not. The ``brown note'' hypothesis is thus physically
grounded---the abdominal resonance is real---but the postulated \emph{airborne}
excitation pathway is energetically negligible.

Future work should validate the boundary condition assumptions through finite
element analysis, extend the model to include gas pocket resonance in the
bowel, and investigate the non-linear regime relevant to high-amplitude
mechanical vibration.
